\documentclass[a4paper,10pt,oneside]{ltjsreport}
\usepackage{fancyhdr}
\usepackage[font=small,labelfont=bf,skip=6pt]{caption}
\usepackage{tcolorbox}
\tcbuselibrary{skins,breakable}
\captionsetup[figure]{format=hang,labelsep=colon}
\linespread{1.15}
\fancyhf{}
\fancyhead[R]{霧窓試験帳}
\fancyfoot[C]{\thepage}
\title{霧窓の採光記録}
\author{薄氷観測室}
\date{}
\tcbset{colback=blue!4!white,colframe=blue!45!black,coltitle=white,fonttitle=\bfseries,boxrule=0.5mm}
\newtcolorbox{observation}[1]{title={#1},sharp corners,left=2mm,right=2mm,top=1mm,bottom=1mm}
\begin{document}
\pagestyle{empty}
\maketitle
\thispagestyle{empty}
\tableofcontents
\clearpage
\pagestyle{fancy}
\chapter{格子戸の測定}
\section{朝の採光}
霧窓を東へ向け、格子の内側に白い紙を置いた。\\
最初の影が印の端に触れるまで、窓の角度を保った。
\section{夕方の採光}
影の伸びる向きだけを紙帯に記す。
\begin{figure}[htbp]
\centering\rule{36mm}{8mm}
\caption[窓の幅を写した帯]{窓の幅を写した帯。\\細い印は開口の端を表す。}\label{fig:band}
\end{figure}
\begin{observation}{凍結時の扱い}
氷の粒は窓枠を外してから数える。
\end{observation}
\newpage
図\ref{fig:band}の紙帯を次の測定にも使う。
\chapter{遮光板の置き方}
\section{二枚の目印}
黒い板を交互に重ね、開口の幅を記録する。
\begin{tcolorbox}[enhanced,breakable,title={保管用の帯},before skip=8pt,after skip=6pt]
乾いた帯の端に風向きを書いた。
\end{tcolorbox}
\end{document}
